\section{An actual Borel order-energy lattice and original current multipliers}
\label{sec:r12-order-energy-lattice}

Retain exactly \textup{(ED1)} and the arbitrary original
middle-positive separately closed background $S$.  The
potential cone remains the frozen
$\mathscr C=\widehat{\mathcal E}^1_{\rm str}(V)$.
We now complete its difference space for the distinct lattice
norm \textup{(FO8)}.  This introduces a larger \emph{scalar
multiplier space}, not a larger potential cone.  Its Borel
realization and the closedness of the old cone are proved
rather than assumed.

\begin{definition}[The order-energy lattice]
\label{def:r12-order-energy-lattice}
Let $\mathbb L_S$ be the norm completion of
$(\mathbb D_S,\|\cdot\|_{\rm o})$.
\end{definition}

The notation $P(f)$ will also be used for a Borel obstacle
$f$ not initially known to lie in $\mathbb D_S$, whenever
at least one member of $\mathscr C$ lies below $f$.
The R11 relative Perron theorem then gives exactly the same
greatest-minorant construction.  This is an explicitly
nonempty-family use of the frozen theorem.

\begin{theorem}[Borel realization and an exact majorant norm]
\label{thm:r12-Borel-order-lattice}
The completion $\mathbb L_S$ is a Banach lattice whose
elements have canonical Borel classes, finite outside a
locally outer $S$-capacity-null set.  These classes are
injectively determined by the completion elements.  Lattice
operations agree with the actual Borel operations, and
\[
 \boxed{\qquad
 \|f\|_{\mathbb L_S}
    =r(P(-|f|))
    =\min\{r(a):a\in\mathscr C,\ a\preceq-|f|\}.
 \qquad}
 \tag{OL1}
\]
In particular every member admits an actual cone majorant
of its absolute value.  The capacity estimate
\textup{(FO12)} holds with this norm.  The original potential
cone $\mathscr C$ is closed in $\mathbb L_S$, and its norm
is still exactly $r$.

More precisely, each $f\in\mathbb L_S$ has a sequence
$f_j\in\mathbb D_S$ and cone functions $W_j\in\mathscr C$
such that
\[
 |f-f_j|\preceq-W_j,\qquad r(W_j)\longrightarrow0.
 \tag{OL2}
\]
The construction does not assert that every Borel function
with some finite-energy cone majorant belongs to $\mathbb L_S$.
\end{theorem}

\noindent\textbf{Status: \texttt{NEW\_PROVED}} as an
identified Borel Banach lattice over the unchanged full
strong potential cone.  The decreasing-series theorem R11
\textup{(CI8)} is an \texttt{INTERNAL\_INPUT}.

\begin{proof}
Take a Cauchy row in the order norm and pass to a subsequence,
still denoted $(f_j)$, with
\[
                 \sum_j\|f_{j+1}-f_j\|_{\rm o}<\infty.
\]
Put $w_j=P(-|f_{j+1}-f_j|)\in\mathscr C$.  Then
$|f_{j+1}-f_j|\preceq-w_j$ and
$\sum_jr(w_j)<\infty$.  The partial sums of the nonpositive
$w_j$ decrease, and Minkowski bounds their energy roots.
R11 \textup{(CI8)} gives their Borel cone sum and every
tail sum
\[
 W_j=\sum_{l\ge j}w_l\in\mathscr C,\qquad
 r(W_j)\le\sum_{l\ge j}r(w_l)\longrightarrow0.
 \tag{OL3}
\]
Outside one countable union of relative null sets, this is
an ordinary series of nonpositive numbers with finite sum.
Thus the series of successive absolute differences is
pointwise convergent there.  It defines a finite Borel
limit $f$ and gives \textup{(OL2)}.  The same argument
identifies the limit in local capacity for the entire
original Cauchy row, by \textup{(FO12)} and the Cauchy
property.

If this limit is the zero Borel class, then
$|f_j|\preceq-W_j$.  The greatest-minorant property gives
$\|f_j\|_{\rm o}\le r(W_j)\to0$.  Thus the completion
element itself is zero.  Applying this observation to
the difference of two Cauchy rows proves both independence
of choices and injectivity of the Borel realization.  This
step excludes an abstract norm vector invisible to Borel
capacity.

Put $a_j=P(-|f_j|)$.  Since
$|f|\preceq-a_j-W_j$, the Borel obstacle $-|f|$ has
a cone minorant $a_j+W_j$.  Its envelope $a=P(-|f|)$
therefore exists, and
\[
 r(a)\le r(a_j+W_j)\le\|f_j\|_{\rm o}+r(W_j).
 \tag{OL4}
\]
Conversely $|f_j|\preceq|f|+|f_j-f|\preceq-a-W_j$,
so
\[
                  \|f_j\|_{\rm o}\le r(a)+r(W_j).
 \tag{OL5}
\]
Passing to the limit proves \textup{(OL1)}.  The minimum
statement follows from energy monotonicity and the greatest
minorant property.  The same majorant gives
\textup{(FO12)} for the completed element.

The pointwise inequalities
$||f|-|g||\le|f-g|$ and the corresponding inequalities
for maxima show that lattice operations are continuous
in the order norm on $\mathbb D_S$.  They extend to the
completion and agree with the Borel limits just constructed.
Formula \textup{(OL1)} also proves directly that the extended
norm is monotone in absolute Borel order.  Thus the completion
is a Banach lattice with the asserted concrete realization.

Finally let a cone row $u_j$ converge in $\mathbb L_S$.
Its roots $r(u_j)=\|u_j\|_{\mathbb L_S}$ are bounded.
The second line of \textup{(FO11)} makes it strongly
Cauchy, and the frozen strong completeness supplies a
limit $u\in\mathscr C$.  Both its strong limit and its
order-norm limit have the same local capacity class.
Injectivity of the Borel realization identifies them.
This proves closedness of $\mathscr C$, without declaring
it equal to every nonpositive member of the larger lattice.
\end{proof}

\begin{proposition}[Relative quasicontinuity in the lattice]
\label{prop:r12-lattice-quasicontinuity}
Every member of $\mathbb L_S$ has a finite representative
which is locally ordinary-open outer $S$-quasicontinuous,
with the same fixed-domain outer capacity as the original
completion.  This does not identify raw and outer capacity
on arbitrary Borel sets.
\end{proposition}

\noindent\textbf{Status: \texttt{DIRECT\_COROLLARY}} of
\textup{(OL2)--(OL3)}, the capacity estimate and the frozen
quasicontinuity of strong cone members.

\begin{proof}
In constructing \textup{(OL2)}, pass to a faster subsequence
so that
$C_0 2^{jd}r(W_j)^d\le2^{-2j}$.
The sets where $|f-f_j|>2^{-j}$ have outer capacity at
most $2^{-2j}$.  Cover a sufficiently late tail of these
sets by ordinary open sets with arbitrarily small total
capacity.  Outside their union the approximants converge
uniformly.  Each $f_j$ is a difference of two strong cone
members and is therefore finite and continuous off a
further ordinary open set of arbitrarily small capacity
on the fixed compact core.  Choose and sum these exceptional
capacities over the countably many approximants.  The uniform
limit is finite and continuous off their union.  A countable
compact exhaustion gives the stated local result.  Only
outer capacity subadditivity and its defining open covers
are used.
\end{proof}

\subsection{The variation norm is the dual norm of this lattice}

\begin{theorem}[Literal curvature integration and weighted total variation]
\label{thm:r12-lattice-dual-variation}
For each $u\in\mathscr C$, the functional $J(u)$ extends
uniquely to $\mathbb L_S$ and satisfies
\[
 J(u)(f)=-\int_V f\,d\widehat\mu_u,\qquad
 \|J(u)\|_{\mathbb L_S^*}=r(u)^q.
 \tag{OL6}
\]
The integral is literal and absolutely convergent for every
completed lattice element.  If $u,v\in\mathscr C$ and
$\sigma=\widehat\mu_u-\widehat\mu_v$, then
\[
 \begin{split}
 \|J(u)-J(v)\|_{\mathbb L_S^*}
   &=\sup_{\substack{a\in\mathscr C\\r(a)\le1}}
                               \int_V(-a)\,d|\sigma|\\
   &=\delta_1(\widehat\mu_u,\widehat\mu_v).
 \end{split}
 \tag{OL7}
\]
Here $\delta_1$ uses the R11 weighted absolute-variation
formula, not the new signed operator distance $\delta_*$
on $\mathbb D_S^*$.  In \textup{(OL7)} the supremum is
unchanged if it is restricted to actual first-energy
potentials, as in the frozen formula.  Applying that
formula to these general completed source measures is
not an import of R11's more restricted dominated-data
stability theorem.
\end{theorem}

\noindent\textbf{Status: \texttt{NEW\_PROVED}} as a
norm-dual interpretation of the two distinct source distances.

\begin{proof}
Since $J(u)$ is order reversing, \textup{(FO10)} gives its
unique continuous extension, with the same norm.  Let
$f\in\mathbb L_S$ and $a=P(-|f|)$.  Then
\[
 \int|f|\,d\widehat\mu_u
       \le\int(-a)\,d\widehat\mu_u
       =B(a,u^{[q]})\le\|f\|_{\mathbb L_S}r(u)^q.
 \tag{OL8}
\]
All relative null sets are ignored by the original curvature
measure.  The approximants in \textup{(OL2)} have integral
error at most $r(W_j)r(u)^q\to0$, proving literal agreement
with the functional extension.  Testing on $u/r(u)$ gives
the norm equality, with the zero case interpreted separately.

The difference $\sigma$ is a locally finite signed
Radon measure; its global total variation may be infinite.
All weighted variation integrals below are finite.
Put $\delta$ equal to the supremum in \textup{(OL7)}.
It is finite, at most $r(u)^q+r(v)^q$.
For any $f\in\mathbb L_S$ of norm at most one, its
minorant $a=P(-|f|)$ has $r(a)\le1$.  Therefore
\[
 \left|\int f\,d\sigma\right|
          \le\int|f|\,d|\sigma|
          \le\int(-a)\,d|\sigma|\le\delta.
 \tag{OL9}
\]
This is one inequality in \textup{(OL7)}.

For the reverse fix $a\in\mathscr C$, $r(a)\le1$.
The finite positive measure $(-a)|\sigma|$ ignores the
joint pole-null set.  By Radon regularity choose disjoint
compact subsets of a Hahn positive and negative carrier
that omit arbitrarily little of this finite weighted measure.
Choose an original compact smooth scalar $\chi$ equal to
$1$ on the first compact and $-1$ on the second, with
$|\chi|\le1$.  The original smooth cutoff and finite
ordinary-open gluing calculus supplies such a function.
For $M>0$ define the lattice clipping
\[
                     f_M=\max\bigl(a,\min(-a,M\chi)\bigr).
 \tag{OL10}
\]
Every entry is in $\mathbb D_S$, including $M\chi$ by
R11 \textup{(HC14)--(HC15)}.  Hence $f_M\in\mathbb D_S$.
It satisfies $|f_M|\preceq-a$, so its order norm is at
most $r(a)\le1$.  On the two chosen compact sets,
$f_M$ tends respectively to $-a$ and $a$ as $M\to\infty$,
outside the harmless pole set.  The error everywhere is
bounded by $2(-a)$, which is integrable for $|\sigma|$.
First let $M\to\infty$ and then let the omitted weighted
mass tend to zero.  The pairings $\int f_M\,d\sigma$
approach $\int(-a)\,d|\sigma|$ from within arbitrarily
small error.  Taking the supremum over $a$ proves the
reverse inequality.  No product $a\chi$ or density of
background currents was required.

Finally, to compare the tester class with the frozen
absolute-variation formula, choose for any $a\in\mathscr C$
an actual strong defining row $a_j$.  The completed tester
estimate \textup{(EV1)} gives
\[
 \int|a_j-a|\,d|\sigma|
  \le\int|a_j-a|\,(d\widehat\mu_u+d\widehat\mu_v)
                                      \longrightarrow0.
\]
Their energy roots converge.  If $r(a)\le1$, replacing
$a_j$ by $a_j/\max(1,r(a_j))$ gives actual unit-root
testers with the same limiting weighted variation
pairing.  When $r(a)=0$, its pairing is zero.  The
actual and completed suprema therefore agree.  This
argument needs neither a common actual lower barrier
nor a common actual dominating source measure.
\end{proof}

\paragraph{Topology qualification.}
The positive functionals $J(u)$ have the same individual
norm in the two dual spaces, but differences of them need
not.  In the original annulus example \textup{(CG1)--(CG16)},
$J(u_s)\to J(u_0)$ in the signed $\mathbb D_S^*$ norm,
while their $\mathbb L_S^*$ distance has lower limit at
least $\sqrt2$.  Thus strong curvature continuity must not
be upgraded to lattice-dual norm continuity.  Moreover,
Theorems~\ref{thm:r12-pole-order-dual-not-strong}
and~\ref{thm:r12-strong-not-order-topology} show that
the positive-dual continuity gate and \textup{(FO13)}
are both \texttt{FALSE}: the order lattice is strictly
stronger than the strong potential topology, and it has
nonzero positive functionals invisible to every original
compact smooth scalar test.  The multiplier conclusions
below need only order-norm convergence and are unaffected.

\subsection{Multipliers of every completed mixed original current}

\begin{lemma}[A finite-energy trace-completing potential]
\label{lem:r12-trace-completing-potential}
Let $0\le k\le q$ and let $\eta$ be an original compact
smooth complementary $(q-k,q-k)$-test form for an original
middle-positive $(p+k,p+k)$-current.  There are a bounded
actual classical $\psi\in\mathscr C$ of finite total top
mass and a finite constant $C_\eta$ such that, for every
tuple $u_1,\ldots,u_k\in\mathscr C$, the coefficient
measure of $T=\widehat N_k(u_1,\ldots,u_k)$ satisfies
\[
 |\mu_{T,\eta}|
        \le C_\eta\,
             \widehat N_q(u_1,\ldots,u_k,\psi^{[q-k]}).
 \tag{OL11}
\]
For $k=q$, one may omit $\psi$ and take
$C_\eta=\|\eta\|_\infty$ for the scalar test.
The domination is local on the coefficient support and
is an inequality of positive Radon measures, not a
closedness assertion for their Borel restrictions.
\end{lemma}

\noindent\textbf{Status: \texttt{DIRECT\_COROLLARY}} of
the frozen original middle-coefficient trace bound, analytic
compensators R11 \textup{(AE2)--(AE5)}, and finite-mass
collar construction \textup{(TC7)--(TC8)}.

\begin{proof}
Use finitely many original moment stages and a compact
smooth partition near the support of $\eta$.  On each
stage a fixed multiple of the full moment trace power
dominates every coefficient of that test in the middle
matrix convention.  The original analytic compensator
lemma gives a global smooth classical psh function whose
curvature dominates this stage's positive $(1,1)$ trace
near the relevant compact piece.  Sum the finitely many
compensators, obtaining one smooth classical $g$ with
all these local domination properties.  Choose $C$ with
$g-C<-1$ on $V$, and choose $A$ large enough that
\[
                         \psi=\max(g-C,A\rho)
\]
equals $g-C$ near the whole test support.  Here $\rho$
is the strict-box barrier.  The exact collar construction
proves that $\psi$ is bounded, actual, zero-boundary in
the frozen energy sense, and has finite total top mass.
Thus it belongs to every positive power class.

On each compact test stage, $D\psi=Dg$ dominates its
full moment trace.  Taking powers preserves the necessary
strong-form domination; the middle-positive coefficient
matrix bound then controls $|\mu_{T,\eta}|$ by a
constant times $(D\psi)^{q-k}\wedge T$ there.
Original smooth-factor compatibility identifies the latter
with the restriction of the mixed top measure in
\textup{(OL11)}.  Sum the finitely many stage bounds.
All moment stages may include any specified actual test
refinement, so this is a complete original-test assertion.
\end{proof}

\begin{theorem}[Order-energy multipliers and closed distributional cap currents]
\label{thm:r12-lattice-current-multipliers}
Let $0\le k\le q$, $u_1,\ldots,u_k\in\mathscr C$, and
$T=\widehat N_k(u_1,\ldots,u_k)$, with $T=S$ when $k=0$.
Every $f\in\mathbb L_S$ defines a literal order-zero
original coefficient product $fT$.  For every compact
original complementary test $\eta$ and the potential
$\psi$ from Lemma~\ref{lem:r12-trace-completing-potential},
\[
 \int_V|f|\,d|\mu_{T,\eta}|
       \le C_\eta\|f\|_{\mathbb L_S}
                   \prod_{i=1}^k r(u_i)\,r(\psi)^{q-k}.
 \tag{OL12}
\]
For $k=q$, the last factor is one.  Thus order-norm
convergence of multipliers gives local coefficient-total-
variation convergence against every original test.

If $k<q$, the formula
\[
                 \mathcal D_T(f)=D(fT)
 \tag{OL13}
\]
defines a linear continuous original current-valued map
whose values are separately $d'$- and $d''$-closed.
For $f\in\mathscr C$ it is exactly the frozen positive
BT current
$\widehat N_{k+1}(f,u_1,\ldots,u_k)$.
For general $f\in\mathbb L_S$, neither positivity nor
order zero of \textup{(OL13)} is asserted.
Every equality and convergence is on the full original
compact smooth complementary test sheaf and every actual
moment refinement.
\end{theorem}

\noindent\textbf{Status: \texttt{NEW\_PROVED}} for the
new completed scalar multiplier domain and its continuous
closed-current action, in every dimension and for arbitrary
original $S$.  The already constructed BT products are
\texttt{INTERNAL\_INPUT}, not newly defined by this theorem.

\begin{proof}
For $f\in\mathbb L_S$ put $a=P(-|f|)$, so
$|f|\preceq-a$ and $r(a)=\|f\|_{\mathbb L_S}$.
The coefficient domination \textup{(OL11)} and the
literal mixed-energy formula give
\[
 \begin{split}
 \int|f|\,d|\mu_{T,\eta}|
 &\le C_\eta\int(-a)\,
            d\widehat N_q(u_1,\ldots,u_k,\psi^{[q-k]})\\
 &=C_\eta B(a,u_1,\ldots,u_k,\psi^{[q-k]})\\
 &\le C_\eta r(a)\prod_i r(u_i)\,r(\psi)^{q-k}.
 \end{split}
\]
The top mixed measure ignores every relative null set;
the coefficient domination therefore shows that the
same is true of the coefficient measure.  Hence the
product is independent of the Borel representative of
$f$.  These coefficient products agree on overlaps by
their ordinary measure definition, and their local
bounds give an original order-zero current $fT$.
Applying the estimate to a difference of multipliers
proves the stated coefficient variation convergence.

Distributional $D$ is continuous on original currents.
It is therefore legitimate to define \textup{(OL13)};
the identities $d'D=0$ and $d''D=0$ give separate
closedness.  At a cone multiplier the frozen original
completed recurrence identifies it with the stated BT
product.  At an exact difference of cone members it is
the corresponding signed difference.  At an arbitrary
lattice limit it is its original current limit, but
no closed positive cone argument has been supplied to
make that distribution a positive order-zero current.
The proof retains this distinction explicitly.
\end{proof}

\paragraph{No relation-quotient shortcut.}
The scalar Borel quotient used in this construction is
already the proved relative potential quotient.  It is
not the unresolved actual analytic PS/PL presentation
quotient.  Extending $f\mapsto D(fT)$ to the lattice
does not show that two actual delta-form presentations
with the same analytic value have the same arbitrary-$S$
action.  Nor does ordinary current equality become an
equality on all additional GK delta tests.  These are
separate gates from the original-test multiplier theorem.
